\BourbakiExercisesSection

\begin{enumerate}
\def\labelenumi{\arabic{enumi}.}
\item
  确定一个具有 \(n\) 个元素的集合上的所有环结构，其中 \(2 \leqslant n \leqslant 7\) ，并确定这些环的理想。
\item
  设 A 是一个环。
\end{enumerate}

\begin{enumerate}
\def\labelenumi{(\alph{enumi})}
\item
  将 \(a \in A\) 对应于左位似 \(x \mapsto ax\) 的映射是 A 到 A 的加法群的与诸右位似交换的自同态环上的一个同构。
\item
  设 M 为 A 的四个元素组成的序列 \((a, b, c, d)\) 的集合，这些序列可以表示成表格或矩阵 \(\begin{pmatrix} a & b \\ c & d \end{pmatrix}\) 的形式。对 M 的每个元素 \(\begin{pmatrix} a & b \\ c & d \end{pmatrix}\) ，我们使映射 \((x, y) \mapsto (ax + by, cx + dy)\) 与之对应，它是 \(\mathbf{A} \times \mathbf{A}\)
\end{enumerate}

到其自身的映射。证明由此得到 M 到与运算 \((x, y) \mapsto (xr, yr)\) （对所有 \(r \in A\)）交换的交换群 A × A 的自同态环 B 上的一个双射。此后将 M 与环 B 等同。计算两个矩阵的乘积。

\begin{enumerate}
\def\labelenumi{(\alph{enumi})}
\setcounter{enumi}{2}
\tightlist
\item
  证明子集 \(\mathbf{T}\) ——由 \(\mathbf{M}\) 中形如 \(\begin{pmatrix} a & b \\ 0 & c \end{pmatrix}\) 的矩阵组成——是 \(\mathbf{M}\) 的一个使子群 \(\mathbf{A} \times 0\) 在 \(\mathbf{A} \times \mathbf{A}\) 中保持稳定的子环。
\end{enumerate}

映射 \(\begin{pmatrix} a & b \\ 0 & c \end{pmatrix} \mapsto a\) 与 \(\begin{pmatrix} a & b \\ 0 & c \end{pmatrix} \mapsto c\) 是 T 到 A 中的同态；设 \(a\) 与 \(b\) 为它们的核。证明 \(a + b = T\) ，\(ab = \{0\}\) ，且 \(ba = b \cap a\) \(\neq \{0\}\) （若 \(A \neq \{0\}\) ；参见第 9 号，命题 6）。

\begin{enumerate}
\def\labelenumi{(\alph{enumi})}
\setcounter{enumi}{3}
\tightlist
\item
  若 \(\mathbf{A} = \mathbf{Z} / 2\mathbf{Z}\) ，则 \(\mathbf{T}\) 不是交换的，但 \(\mathbf{T}\) 的可逆元素的乘法群是交换的。
\end{enumerate}

\begin{enumerate}
\def\labelenumi{\arabic{enumi}.}
\setcounter{enumi}{2}
\item
  设 A 为一个环，且 \(a \in A\) 。如果存在且仅存在一个 \(a' \in A\) 使得 \(aa' = 1\) ，则 a 是可逆的，并且 \(a' = a^{-1}\) （首先证明 a 是左可消去的，然后考虑乘积 \(aa'a\)）。
\item
  设 A 为一个伪环，其中 A 的每个加法子群都是 A 的左理想。
\end{enumerate}

\begin{enumerate}
\def\labelenumi{(\alph{enumi})}
\item
  若 \(\mathbf{A}\) 是一个环，则 \(\mathbf{A}\) 典范同构于 \(\mathbf{Z} / n\mathbf{Z}\) ，其中 \(n\) 为某个适当的整数。
\item
  如果 A 没有单位元且每个非零元素都是可消去的，则 A 同构于 Z 的一个伪子环。
\end{enumerate}

\begin{enumerate}
\def\labelenumi{\arabic{enumi}.}
\setcounter{enumi}{4}
\tightlist
\item
  设 \(\mathbf{M}\) 为一个幺半群，\(\mathbf{Z}^{(\mathbf{M})}\) 为以 \(\mathbf{M}\) 为基的自由交换群，\(u\colon \mathbf{M}\to \mathbf{Z}^{(\mathbf{M})}\) 为典范映射。
\end{enumerate}

\begin{enumerate}
\def\labelenumi{(\alph{enumi})}
\item
  证明在 \(\mathbf{Z}^{(\mathbf{M})}\) 上存在唯一的乘法，使得 \(\mathbf{Z}^{(\mathbf{M})}\) 是一个环，且 \(u\) 是幺半群 \(\mathbf{M}\) 到环 \(\mathbf{Z}^{(\mathbf{M})}\) 的乘法幺半群中的幺半群态射。
\item
  设 A 为一个环，\(v: \mathbf{M} \to \mathbf{A}\) 为幺半群 M 到 A 的乘法幺半群中的幺半群态射。存在唯一的环态射 \(f: \mathbf{Z}^{(\mathbf{M})} \to \mathbf{A}\) ，使得
\end{enumerate}

\[
f \circ u = v.
\]

\begin{enumerate}
\def\labelenumi{\arabic{enumi}.}
\setcounter{enumi}{5}
\tightlist
\item
  设 A 为一个交换环，M 为 A 的加法群的一个子幺半群。设 n 为一个整数 \textgreater0 ，使得 \(m^{n} = 0\) 对 \(m \in M\) 成立。
\end{enumerate}

\begin{enumerate}
\def\labelenumi{(\alph{enumi})}
\item
  若 \(n!\) 在 \(\mathbf{A}\) 中不是零因子，则由 \(n\) 个 \(\mathbf{M}\) 中元素组成的族之乘积为零。
\item
  证明：若 \(n!\) 在 A 中是零因子，则 (a) 中的结论不一定成立。
\end{enumerate}

\begin{enumerate}
\def\labelenumi{\arabic{enumi}.}
\setcounter{enumi}{6}
\tightlist
\item
  对于所有 \(n \in \mathbf{N}^*\) ,
\end{enumerate}

\[
n! = \sum_ {\nu = 0} ^ {n - 1} (- 1) ^ {\nu} \binom {n} {\nu} (n - \nu) ^ {n}
\]

（利用第 2 号，命题 2。）

\begin{enumerate}
\def\labelenumi{\arabic{enumi}.}
\setcounter{enumi}{7}
\item
  在环中，由一个左理想生成的右理想是一个双边理想。
\item
  在环中，一个右理想的右零化子是一个双边理想。
\item
  在环 A 中，由元素 xy - yx（其中 x 与 y 遍历 A）生成的双边理想，是使 A/a 为交换环的双边理想 a 中最小的一个。
\end{enumerate}

¶ 11. 在环 A 中，一个双边理想 a 称为不可约的，如果不存在有序对 (b, c)，其中 b、c 是与 a 不同的双边理想，且满足 a = b ∩ c。

\begin{enumerate}
\def\labelenumi{(\alph{enumi})}
\item
  证明 A 的所有不可约双边理想的交化为 0（注意，不含任何元素 \(a \neq 0\) 的双边理想组成的集合是归纳的，并应用佐恩引理）。
\item
  由此推出，每个双边理想 \(\mathfrak{a}\) （\(\mathbf{A}\) 的）是包含它的所有不可约理想的交。
\end{enumerate}

\begin{enumerate}
\def\labelenumi{\arabic{enumi}.}
\setcounter{enumi}{11}
\tightlist
\item
  设 A 为一个环，且 \((\mathfrak{a}_{i})_{i\in I}\) 为 A 的左理想的有限族，使得 A 的加法群是诸加法群 \(a_{i}\) 的直和。若 \(1=\sum_{i\in I}e_{i}(e_{i}\in\mathfrak{a}_{i})\) ，则 \(e_{i}^{2}=e_{i}, e_{i}e_{j}=0\) （若 \(i\neq j\)）且 \(a_{i}=Ae_{i}\) （对所有 \(x\in A\) 写出 x=x.1）。反之，如果 \((e_{i})_{i\in I}\) 是 A 的幂等元的有限族，满足 \(e_{i}e_{j}=0\) （对 \(i\neq j\)）且 \(1=\sum_{i}e_{i}\) ，则 A 是诸左理想 \(Ae_{i}\) 的直和。为使诸理想 \(Ae_{i}\) 是双边的，必要且充分条件是诸 \(e_{i}\) 属于 A 的中心。
\end{enumerate}

¶ 13. 设 A 是一个环，e 是 A 的一个幂等元。

\begin{enumerate}
\def\labelenumi{(\alph{enumi})}
\item
  证明 \(\mathbf{A}\) 的加法群是左理想 \(\mathfrak{a} = \mathbf{A}e\) 与左零化子 \(\mathfrak{b}\) （对 \(e\) 而言）的直和（注意，对所有 \(x\in \mathbf{A}\) ，有 \(x - xe\in \mathfrak{b}\)）。
\item
  每个右理想 \(\mathfrak{d}\) （属于 \(\mathbf{A}\)）是 \(\mathfrak{d} \cap \mathfrak{a}\) 与 \(\mathfrak{d} \cap \mathfrak{b}\) 的直和。
\item
  若 \(\mathbf{Ae} = e\mathbf{A}\) ，则 \(\mathfrak{a}\) 与 \(\mathfrak{b}\) 是 \(\mathbf{A}\) 的双边理想，并且定义 \(\mathbf{A}\) 的一个直分解。
\end{enumerate}

\begin{enumerate}
\def\labelenumi{\arabic{enumi}.}
\setcounter{enumi}{13}
\item
  设 A 为一个交换环，其中零因子只有有限 n 个。证明 A 至多有 \((n+1)^{2}\) 个元素。（设 \(a_{1},\ldots,a_{n}\) 为零因子。证明零化子 \(J_{1}\) （\(a_{1}\) 的）至多有 \(n+1\) 个元素，并且商环 A/ \(J_{1}\) 至多有 \(n+1\) 个元素。）
\item
  环胚是一个带有两个合成律的集合 E：（a）一个结合乘法 \(xy\) ；（b）一个以加法记、并非处处有定义的合成律（§ 1，第 1 号），且满足下列条件：
\end{enumerate}

\begin{enumerate}
\def\labelenumi{(\arabic{enumi})}
\item
  它是交换的（换句话说，若 \(x + y\) 有定义，则 \(y + x\) 亦有定义且 \(x + y = y + x\) ；此时称 \(x\) 与 \(y\) 是可相加的）；
\item
  若 \(x\) 与 \(y\) 可相加，则为使 \(x + y\) 与 \(z\) 可相加，只需 \(x\) 与 \(y\) 、以及 \(y\) 与 \(z\) 分别可相加；此时 \(x\) 与 \(y + z\) 可相加，且 \((x + y) + z = x + (y + z)\) ；
\item
  存在一个单位元 0；
\item
  若 \(x\) 与 \(z\) 、以及 \(y\) 与 \(z\) 分别可相加，且 \(x + z = y + z\) ，则 \(x = y\) ；
\item
  若 \(x, y\) 可相加，则 \(xz\) 与 \(yz\) （相应地 \(zx\) 与 \(zy\)）亦可相加，且
\end{enumerate}

\[
(x + y) z = x z + y z
\]

（相应地 \(z(x + y) = zx + zy)\) 对所有 \(z \in \mathbf{E}\) 成立。

每个环都是一个环状结构。

\begin{enumerate}
\def\labelenumi{(\alph{enumi})}
\item
  考察 § 8 的定义和结果以及上述习题如何推广到环胚（环胚 E 的左理想是 E 的一个子集 \(\mathfrak{a}\) ，它在加法下稳定，且满足 \(\mathbf{E}.\mathfrak{a}\subset \mathfrak{a}\)）。
\item
  设 G 为一个带算子的群，f 与 g 为 G 的两个自同态。为使映射 \(x \mapsto f(x)g(x)\) 是 G 的自同态，必要且充分条件是子群 \(f(\mathrm{G})\) 的每个元素与 \(g(\mathrm{G})\) 的每个元素可交换；若把这一自同态记作 \(f + g\) ，且 fg 为复合自同态 \(x \mapsto f(g(x))\) ，证明 G 的具有这些合成律的自同态集合 E 是一个有单位元的环胚；为使 E 成为一个环，必要且充分条件是 G 是交换的。
\end{enumerate}

一个元素 \(f \in E\) 能与 E 的所有元素相加，必要且充分条件是 \(f(G)\) 含于 G 的中心；这些自同态组成的集合 N 是一个伪环，称为环胚 E 的核。

一个自同态 \(f\) （\(G\) 的）称为正规的，如果它与 \(G\) 的所有内自同构可交换；对每个正规稳定子群 \(H\) ，\(G, f(G)\) 都是 \(G\) 的一个正规稳定子群。证明 \(D\) ——\(G\) 的正规自同态组成的集合——构成 \(E\) 的一个子环胚，且核 \(N\) 是一个双边理想，位于 \(D.†\)

\begin{enumerate}
\def\labelenumi{\arabic{enumi}.}
\setcounter{enumi}{15}
\tightlist
\item
  给出理想 \(a, b, c\) （在环 \(\mathbf{Z}\) 中）的例子，使得 \(ab \neq a \cap b\) 且 \((a + b)(a + c) \neq a + bc\) 。
\end{enumerate}

